%% ===================================================
%% ALESSIO ROMA — Curriculum Artistico
%% Compilare con: xelatex curriculum-artistico.tex
%% Versione chiara (sfondo bianco):
%%   xelatex -jobname=curriculum-artistico-chiaro "\def\lightmode{}%% ===================================================
%% ALESSIO ROMA — Curriculum Artistico
%% Compilare con: xelatex curriculum-artistico.tex
%% Versione chiara (sfondo bianco):
%%   xelatex -jobname=curriculum-artistico-chiaro "\def\lightmode{}%% ===================================================
%% ALESSIO ROMA — Curriculum Artistico
%% Compilare con: xelatex curriculum-artistico.tex
%% Versione chiara (sfondo bianco):
%%   xelatex -jobname=curriculum-artistico-chiaro "\def\lightmode{}%% ===================================================
%% ALESSIO ROMA — Curriculum Artistico
%% Compilare con: xelatex curriculum-artistico.tex
%% Versione chiara (sfondo bianco):
%%   xelatex -jobname=curriculum-artistico-chiaro "\def\lightmode{}\input{curriculum-artistico}"
%% ===================================================
\documentclass[a4paper, 10pt]{article}

\usepackage{fontspec}
\usepackage{xcolor}
\usepackage{geometry}
\usepackage{graphicx}
\usepackage{hyperref}
\usepackage{enumitem}
\usepackage{multicol}
\usepackage{setspace}

%% ——— Font ————————————————————————————————————
\setmainfont[
  Path           = assets/fonts/,
  UprightFont    = Champagne___Limousines.ttf,
  BoldFont       = Champagne___Limousines_Bold.ttf,
  ItalicFont     = Champagne___Limousines_Italic.ttf,
  BoldItalicFont = Champagne___Limousines_Bold_Italic.ttf
]{Champagne___Limousines.ttf}

%% ——— Colori dal sito ————————————————————————
\ifdefined\lightmode
  %% Versione chiara: stessi toni freddi del sito, scuriti per la stampa
  \definecolor{colbg}{HTML}{FFFFFF}
  \definecolor{coltextmain}{HTML}{1B2124}
  \definecolor{coltextlight}{HTML}{3F4C53}
  \definecolor{colaccent}{HTML}{2A5670}
  \definecolor{colborder}{HTML}{B9C4CA}
\else
  \definecolor{colbg}{HTML}{0D0F10}
  \definecolor{coltextmain}{HTML}{E8E0D0}
  \definecolor{coltextlight}{HTML}{7A8A90}
  \definecolor{colaccent}{HTML}{B8C8D4}
  \definecolor{colborder}{HTML}{2E3A42}
\fi

%% ——— Impostazioni pagina ————————————————————
\geometry{a4paper, left=2.2cm, right=2.2cm, top=2cm, bottom=2cm}
\pagecolor{colbg}
\color{coltextmain}
\setlength{\parindent}{0pt}
\setlength{\parskip}{0pt}

%% ——— Link ————————————————————————————————————
\hypersetup{
  colorlinks = true,
  urlcolor   = colaccent,
  linkcolor  = colaccent,
  pdfauthor  = {Alessio Roma},
  pdftitle   = {Alessio Roma — Curriculum Artistico}
}

%% ——— Helper: testo con spaziatura lettere ————
\newcommand{\spaced}[2][80]{%
  {\addfontfeatures{LetterSpace=#1}#2}%
}

%% ——— Linea divisoria ————————————————————————
\newcommand{\cvrule}{%
  {\color{colborder}\rule{\linewidth}{0.4pt}}%
}

%% ——— Intestazione sezione ————————————————————
\newcommand{\cvsection}[1]{%
  \vspace{0.7em}%
  {\small\bfseries\spaced[90]{\MakeUppercase{#1}}}%
  \vspace{0.15em}\newline
  \cvrule%
  \vspace{0.4em}\newline
}

%% ——— Voce opera: anno, titolo (italic), organico ——
\newcommand{\workrow}[3]{%
  {\footnotesize\color{coltextlight}\spaced[90]{#1}}\enspace
  {\small\itshape #2}\newline
  {\footnotesize\color{coltextlight}\hspace{2.2em}#3}%
  \vspace{0.1em}\newline
}

%% ——— Voce formazione / premio ————————————————
\newcommand{\cvrow}[3]{%
  {\small\bfseries #1}\hfill{\footnotesize\color{coltextlight}\spaced[90]{#2}}%
  \newline
  {\small\color{coltextlight}#3}%
  \vspace{0.55em}\newline
}

%% ===================================================
\begin{document}
\pagestyle{empty}

%% ——— INTESTAZIONE ———————————————————————————

\noindent\begin{minipage}[t]{0.60\linewidth}
  \vspace{0pt}
  {\fontsize{34}{38}\selectfont\spaced[30]{ALESSIO ROMA}}%
  \vspace{0.3em}\newline
  {\footnotesize\bfseries\color{coltextlight}\spaced[130]{MUSICISTA · COMPOSITORE}}%
  \vspace{0.55em}\newline
  {\color{colborder}\rule{0.97\linewidth}{0.4pt}}%
  \vspace{0.45em}\newline
  {\small\color{coltextlight}%
    Bari, 1993\enspace·\enspace
    \href{mailto:info@alessioroma.it}{info@alessioroma.it}\enspace·\enspace
    \href{https://alessioroma.it}{alessioroma.it}%
  }%
\end{minipage}%
\hfill%
\begin{minipage}[t]{0.31\linewidth}
  \vspace{0pt}
  \raggedleft
  \begingroup
    \setlength{\fboxsep}{0pt}%
    \setlength{\fboxrule}{1pt}%
    \color{coltextmain}%
    \fbox{%
      \includegraphics[width=3.8cm, height=3.8cm]%
        {assets/images/alessio-roma-smile-crop.jpg}%
    }%
  \endgroup
\end{minipage}

\vspace{0.9em}
\cvrule

%% ——— PROFILO ————————————————————————————————
\vspace{0.7em}
{\small Alessio Roma è un musicista e compositore nato a Bari nel 1993.
Nel suo lavoro integra elementi extra-musicali — salute mentale,
oppressione dei popoli, violenza sulle donne — attraverso le tecniche
spettraliste e la ricerca sulla fusione timbrica di più strumenti in
un unico timbro percepibile. Le sue opere sono eseguite in Italia e
all'estero in concerti, festival e rassegne di musica contemporanea.}

%% ——— FORMAZIONE ————————————————————————————
\cvsection{Formazione}

\cvrow{Conservatorio «Niccolò Piccinni», Bari}{2026}{%
  Diploma accademico di II livello in Composizione%
  \enspace·\enspace docente M°\,Biagio Putignano}

%% ——— MASTERCLASS ————————————————————————————
\cvsection{Masterclass internazionali}

{\small
Alessandro Melchiorre\enspace·\enspace
Ivan Fedele\enspace·\enspace
Giorgio Nottoli\enspace·\enspace
Marco Stroppa\enspace·\enspace
Robert HP Platz\enspace·\enspace
Roberto Paci Dalò\enspace·\enspace
Juha T.\,Koskinen\enspace·\enspace
Musikschule der Stadt Freising, Germania%
}

%% ——— PREMI ——————————————————————————————————
\cvsection{Premi e riconoscimenti}

{\small 2° Premio, Concorso di composizione «Note per una stella», Monteroni di Lecce (LE)}%
\newline
{\small\color{coltextlight}\hspace{0.5em}per \textit{Cinque piccole variazioni per tromba}}%
\vspace{0.55em}\newline

{\small Due prime esecuzioni assolute, eventi del centenario della\\
Società Italiana di Musica Contemporanea (SIMC)}%
\vspace{0.55em}\newline

%% ——— COMMISSIONI ————————————————————————————
\cvsection{Commissioni}

{\footnotesize\color{coltextlight}\spaced[90]{2026}}\enspace
{\small\textit{Ostinate Migrazioni} per quartetto d'archi — Traetta Opera Festival}%
\vspace{0.55em}\newline

{\footnotesize\color{coltextlight}\spaced[90]{2025}}\enspace
{\small\textit{Chiari Ombre} per violino e violoncello — Traetta Opera Festival}%
\vspace{0.55em}\newline

%% ——— COLLABORAZIONI ——————————————————————————
\cvsection{Collaborazioni}

{\small LEAP — Laboratorio ElettroAcustico Permanente, Roma}%
\newline
{\small\color{coltextlight}\hspace{0.5em}con Giuseppe Silvi (cofondatore)}

%% ——— DISCOGRAFIA ————————————————————————————
\cvsection{Discografia}

{\small Registrazioni disponibili su Spotify, Apple Music, YouTube Music e altre
piattaforme di streaming musicale.}

%% ===================================================
%% PAGINA 2 — CATALOGO DELLE OPERE
%% ===================================================
\newpage

\noindent
{\footnotesize\color{coltextlight}\spaced[100]{ALESSIO ROMA}}%
\hfill%
{\footnotesize\color{coltextlight}\spaced[100]{CATALOGO DELLE OPERE}}%
\vspace{0.3em}\newline
\cvrule

\setlength{\columnsep}{1.8em}
\begin{multicols}{2}

\cvsection{Musica strumentale}

\workrow{2026}{Ostinate Migrazioni}%
  {quartetto d'archi%
  \newline\hspace*{2.2em}{\itshape commissione Traetta Opera Festival}}

\workrow{2026}{Io sono come il mare}%
  {ensemble, elett. dal vivo, voce narrante — 45' ca.%
  \newline\hspace*{2.2em}testi di Antonella Diacono}

\workrow{2025}{Chiari Ombre}%
  {violino e violoncello — 6' ca.%
  \newline\hspace*{2.2em}{\itshape commissione Traetta Opera Festival}}

\workrow{2025}{Torsioni del reale}%
  {violino solo — 7' ca.}

\workrow{2025}{Fecero un deserto e lo chiamarono Pace}%
  {ensemble, elett. dal vivo, voce narrante — 15' ca.}

\workrow{2024}{Serenata per un attacco di panico III}%
  {pianoforte solo — 10' ca.}

\workrow{2024}{Costellazioni}%
  {orchestra — 9' ca.}

\workrow{2024}{Buttate in cielo}%
  {ensemble e narratore — 30' ca.%
  \newline\hspace*{2.2em}testi di Denise De Scisciolo}

\workrow{2023}{Paysage}%
  {organo storico e due performer — 8' ca.}

\workrow{2023}{Posters}%
  {clarinetto e pianoforte — 9' ca.}

\workrow{2023}{Serenata per un attacco di panico II}%
  {sassofono contralto solo — 8' ca.}

\workrow{2023}{Parla, inciampa, rifletti e splendi!}%
  {violino solo — 10' ca.}

\workrow{2023}{It's a long, long way to therapy}%
  {quartetto d'archi e elett. dal vivo — 8' ca.}

\workrow{2023{\addfontfeatures{LetterSpace=0}\;{\tiny(rev.\,2024)}}}{Serenata per un attacco di panico I}%
  {viola sola — 10' ca.}

\workrow{2022}{Fili d'aquilone}%
  {violino, violoncello e pianoforte — 10' ca.}

\workrow{2022}{Il Gioco Grande}%
  {sax baritono, percussioni e narratori — 15'–18' ca.}

\workrow{2022}{Boards}%
  {pianoforte solo — 7' ca.}

\workrow{2021}{Umarmung}%
  {duo di flauto dolce — 6' ca.}

\cvsection{Colonne sonore}

\workrow{2023}{Nu ffischia}%
  {cortometraggio di Minafra, Carella e Piccinni}

\workrow{2021}{Myosotys}%
  {cortometraggio di Minafra e Piccinni}

\workrow{2021}{Punto e a caso}%
  {cortometraggio di Minafra e Piccinni}

\end{multicols}

%% ——— Footer pagina 2 ————————————————————————
\vfill
\cvrule
\vspace{0.35em}
{\footnotesize\color{coltextlight}%
  \href{https://alessioroma.it}{alessioroma.it}%
  \enspace·\enspace
  \href{mailto:info@alessioroma.it}{info@alessioroma.it}%
  \enspace·\enspace
  Aggiornato: ottobre 2026%
}

\end{document}
"
%% ===================================================
\documentclass[a4paper, 10pt]{article}

\usepackage{fontspec}
\usepackage{xcolor}
\usepackage{geometry}
\usepackage{graphicx}
\usepackage{hyperref}
\usepackage{enumitem}
\usepackage{multicol}
\usepackage{setspace}

%% ——— Font ————————————————————————————————————
\setmainfont[
  Path           = assets/fonts/,
  UprightFont    = Champagne___Limousines.ttf,
  BoldFont       = Champagne___Limousines_Bold.ttf,
  ItalicFont     = Champagne___Limousines_Italic.ttf,
  BoldItalicFont = Champagne___Limousines_Bold_Italic.ttf
]{Champagne___Limousines.ttf}

%% ——— Colori dal sito ————————————————————————
\ifdefined\lightmode
  %% Versione chiara: stessi toni freddi del sito, scuriti per la stampa
  \definecolor{colbg}{HTML}{FFFFFF}
  \definecolor{coltextmain}{HTML}{1B2124}
  \definecolor{coltextlight}{HTML}{3F4C53}
  \definecolor{colaccent}{HTML}{2A5670}
  \definecolor{colborder}{HTML}{B9C4CA}
\else
  \definecolor{colbg}{HTML}{0D0F10}
  \definecolor{coltextmain}{HTML}{E8E0D0}
  \definecolor{coltextlight}{HTML}{7A8A90}
  \definecolor{colaccent}{HTML}{B8C8D4}
  \definecolor{colborder}{HTML}{2E3A42}
\fi

%% ——— Impostazioni pagina ————————————————————
\geometry{a4paper, left=2.2cm, right=2.2cm, top=2cm, bottom=2cm}
\pagecolor{colbg}
\color{coltextmain}
\setlength{\parindent}{0pt}
\setlength{\parskip}{0pt}

%% ——— Link ————————————————————————————————————
\hypersetup{
  colorlinks = true,
  urlcolor   = colaccent,
  linkcolor  = colaccent,
  pdfauthor  = {Alessio Roma},
  pdftitle   = {Alessio Roma — Curriculum Artistico}
}

%% ——— Helper: testo con spaziatura lettere ————
\newcommand{\spaced}[2][80]{%
  {\addfontfeatures{LetterSpace=#1}#2}%
}

%% ——— Linea divisoria ————————————————————————
\newcommand{\cvrule}{%
  {\color{colborder}\rule{\linewidth}{0.4pt}}%
}

%% ——— Intestazione sezione ————————————————————
\newcommand{\cvsection}[1]{%
  \vspace{0.7em}%
  {\small\bfseries\spaced[90]{\MakeUppercase{#1}}}%
  \vspace{0.15em}\newline
  \cvrule%
  \vspace{0.4em}\newline
}

%% ——— Voce opera: anno, titolo (italic), organico ——
\newcommand{\workrow}[3]{%
  {\footnotesize\color{coltextlight}\spaced[90]{#1}}\enspace
  {\small\itshape #2}\newline
  {\footnotesize\color{coltextlight}\hspace{2.2em}#3}%
  \vspace{0.1em}\newline
}

%% ——— Voce formazione / premio ————————————————
\newcommand{\cvrow}[3]{%
  {\small\bfseries #1}\hfill{\footnotesize\color{coltextlight}\spaced[90]{#2}}%
  \newline
  {\small\color{coltextlight}#3}%
  \vspace{0.55em}\newline
}

%% ===================================================
\begin{document}
\pagestyle{empty}

%% ——— INTESTAZIONE ———————————————————————————

\noindent\begin{minipage}[t]{0.60\linewidth}
  \vspace{0pt}
  {\fontsize{34}{38}\selectfont\spaced[30]{ALESSIO ROMA}}%
  \vspace{0.3em}\newline
  {\footnotesize\bfseries\color{coltextlight}\spaced[130]{MUSICISTA · COMPOSITORE}}%
  \vspace{0.55em}\newline
  {\color{colborder}\rule{0.97\linewidth}{0.4pt}}%
  \vspace{0.45em}\newline
  {\small\color{coltextlight}%
    Bari, 1993\enspace·\enspace
    \href{mailto:info@alessioroma.it}{info@alessioroma.it}\enspace·\enspace
    \href{https://alessioroma.it}{alessioroma.it}%
  }%
\end{minipage}%
\hfill%
\begin{minipage}[t]{0.31\linewidth}
  \vspace{0pt}
  \raggedleft
  \begingroup
    \setlength{\fboxsep}{0pt}%
    \setlength{\fboxrule}{1pt}%
    \color{coltextmain}%
    \fbox{%
      \includegraphics[width=3.8cm, height=3.8cm]%
        {assets/images/alessio-roma-smile-crop.jpg}%
    }%
  \endgroup
\end{minipage}

\vspace{0.9em}
\cvrule

%% ——— PROFILO ————————————————————————————————
\vspace{0.7em}
{\small Alessio Roma è un musicista e compositore nato a Bari nel 1993.
Nel suo lavoro integra elementi extra-musicali — salute mentale,
oppressione dei popoli, violenza sulle donne — attraverso le tecniche
spettraliste e la ricerca sulla fusione timbrica di più strumenti in
un unico timbro percepibile. Le sue opere sono eseguite in Italia e
all'estero in concerti, festival e rassegne di musica contemporanea.}

%% ——— FORMAZIONE ————————————————————————————
\cvsection{Formazione}

\cvrow{Conservatorio «Niccolò Piccinni», Bari}{2026}{%
  Diploma accademico di II livello in Composizione%
  \enspace·\enspace docente M°\,Biagio Putignano}

%% ——— MASTERCLASS ————————————————————————————
\cvsection{Masterclass internazionali}

{\small
Alessandro Melchiorre\enspace·\enspace
Ivan Fedele\enspace·\enspace
Giorgio Nottoli\enspace·\enspace
Marco Stroppa\enspace·\enspace
Robert HP Platz\enspace·\enspace
Roberto Paci Dalò\enspace·\enspace
Juha T.\,Koskinen\enspace·\enspace
Musikschule der Stadt Freising, Germania%
}

%% ——— PREMI ——————————————————————————————————
\cvsection{Premi e riconoscimenti}

{\small 2° Premio, Concorso di composizione «Note per una stella», Monteroni di Lecce (LE)}%
\newline
{\small\color{coltextlight}\hspace{0.5em}per \textit{Cinque piccole variazioni per tromba}}%
\vspace{0.55em}\newline

{\small Due prime esecuzioni assolute, eventi del centenario della\\
Società Italiana di Musica Contemporanea (SIMC)}%
\vspace{0.55em}\newline

%% ——— COMMISSIONI ————————————————————————————
\cvsection{Commissioni}

{\footnotesize\color{coltextlight}\spaced[90]{2026}}\enspace
{\small\textit{Ostinate Migrazioni} per quartetto d'archi — Traetta Opera Festival}%
\vspace{0.55em}\newline

{\footnotesize\color{coltextlight}\spaced[90]{2025}}\enspace
{\small\textit{Chiari Ombre} per violino e violoncello — Traetta Opera Festival}%
\vspace{0.55em}\newline

%% ——— COLLABORAZIONI ——————————————————————————
\cvsection{Collaborazioni}

{\small LEAP — Laboratorio ElettroAcustico Permanente, Roma}%
\newline
{\small\color{coltextlight}\hspace{0.5em}con Giuseppe Silvi (cofondatore)}

%% ——— DISCOGRAFIA ————————————————————————————
\cvsection{Discografia}

{\small Registrazioni disponibili su Spotify, Apple Music, YouTube Music e altre
piattaforme di streaming musicale.}

%% ===================================================
%% PAGINA 2 — CATALOGO DELLE OPERE
%% ===================================================
\newpage

\noindent
{\footnotesize\color{coltextlight}\spaced[100]{ALESSIO ROMA}}%
\hfill%
{\footnotesize\color{coltextlight}\spaced[100]{CATALOGO DELLE OPERE}}%
\vspace{0.3em}\newline
\cvrule

\setlength{\columnsep}{1.8em}
\begin{multicols}{2}

\cvsection{Musica strumentale}

\workrow{2026}{Ostinate Migrazioni}%
  {quartetto d'archi%
  \newline\hspace*{2.2em}{\itshape commissione Traetta Opera Festival}}

\workrow{2026}{Io sono come il mare}%
  {ensemble, elett. dal vivo, voce narrante — 45' ca.%
  \newline\hspace*{2.2em}testi di Antonella Diacono}

\workrow{2025}{Chiari Ombre}%
  {violino e violoncello — 6' ca.%
  \newline\hspace*{2.2em}{\itshape commissione Traetta Opera Festival}}

\workrow{2025}{Torsioni del reale}%
  {violino solo — 7' ca.}

\workrow{2025}{Fecero un deserto e lo chiamarono Pace}%
  {ensemble, elett. dal vivo, voce narrante — 15' ca.}

\workrow{2024}{Serenata per un attacco di panico III}%
  {pianoforte solo — 10' ca.}

\workrow{2024}{Costellazioni}%
  {orchestra — 9' ca.}

\workrow{2024}{Buttate in cielo}%
  {ensemble e narratore — 30' ca.%
  \newline\hspace*{2.2em}testi di Denise De Scisciolo}

\workrow{2023}{Paysage}%
  {organo storico e due performer — 8' ca.}

\workrow{2023}{Posters}%
  {clarinetto e pianoforte — 9' ca.}

\workrow{2023}{Serenata per un attacco di panico II}%
  {sassofono contralto solo — 8' ca.}

\workrow{2023}{Parla, inciampa, rifletti e splendi!}%
  {violino solo — 10' ca.}

\workrow{2023}{It's a long, long way to therapy}%
  {quartetto d'archi e elett. dal vivo — 8' ca.}

\workrow{2023{\addfontfeatures{LetterSpace=0}\;{\tiny(rev.\,2024)}}}{Serenata per un attacco di panico I}%
  {viola sola — 10' ca.}

\workrow{2022}{Fili d'aquilone}%
  {violino, violoncello e pianoforte — 10' ca.}

\workrow{2022}{Il Gioco Grande}%
  {sax baritono, percussioni e narratori — 15'–18' ca.}

\workrow{2022}{Boards}%
  {pianoforte solo — 7' ca.}

\workrow{2021}{Umarmung}%
  {duo di flauto dolce — 6' ca.}

\cvsection{Colonne sonore}

\workrow{2023}{Nu ffischia}%
  {cortometraggio di Minafra, Carella e Piccinni}

\workrow{2021}{Myosotys}%
  {cortometraggio di Minafra e Piccinni}

\workrow{2021}{Punto e a caso}%
  {cortometraggio di Minafra e Piccinni}

\end{multicols}

%% ——— Footer pagina 2 ————————————————————————
\vfill
\cvrule
\vspace{0.35em}
{\footnotesize\color{coltextlight}%
  \href{https://alessioroma.it}{alessioroma.it}%
  \enspace·\enspace
  \href{mailto:info@alessioroma.it}{info@alessioroma.it}%
  \enspace·\enspace
  Aggiornato: ottobre 2026%
}

\end{document}
"
%% ===================================================
\documentclass[a4paper, 10pt]{article}

\usepackage{fontspec}
\usepackage{xcolor}
\usepackage{geometry}
\usepackage{graphicx}
\usepackage{hyperref}
\usepackage{enumitem}
\usepackage{multicol}
\usepackage{setspace}

%% ——— Font ————————————————————————————————————
\setmainfont[
  Path           = assets/fonts/,
  UprightFont    = Champagne___Limousines.ttf,
  BoldFont       = Champagne___Limousines_Bold.ttf,
  ItalicFont     = Champagne___Limousines_Italic.ttf,
  BoldItalicFont = Champagne___Limousines_Bold_Italic.ttf
]{Champagne___Limousines.ttf}

%% ——— Colori dal sito ————————————————————————
\ifdefined\lightmode
  %% Versione chiara: stessi toni freddi del sito, scuriti per la stampa
  \definecolor{colbg}{HTML}{FFFFFF}
  \definecolor{coltextmain}{HTML}{1B2124}
  \definecolor{coltextlight}{HTML}{3F4C53}
  \definecolor{colaccent}{HTML}{2A5670}
  \definecolor{colborder}{HTML}{B9C4CA}
\else
  \definecolor{colbg}{HTML}{0D0F10}
  \definecolor{coltextmain}{HTML}{E8E0D0}
  \definecolor{coltextlight}{HTML}{7A8A90}
  \definecolor{colaccent}{HTML}{B8C8D4}
  \definecolor{colborder}{HTML}{2E3A42}
\fi

%% ——— Impostazioni pagina ————————————————————
\geometry{a4paper, left=2.2cm, right=2.2cm, top=2cm, bottom=2cm}
\pagecolor{colbg}
\color{coltextmain}
\setlength{\parindent}{0pt}
\setlength{\parskip}{0pt}

%% ——— Link ————————————————————————————————————
\hypersetup{
  colorlinks = true,
  urlcolor   = colaccent,
  linkcolor  = colaccent,
  pdfauthor  = {Alessio Roma},
  pdftitle   = {Alessio Roma — Curriculum Artistico}
}

%% ——— Helper: testo con spaziatura lettere ————
\newcommand{\spaced}[2][80]{%
  {\addfontfeatures{LetterSpace=#1}#2}%
}

%% ——— Linea divisoria ————————————————————————
\newcommand{\cvrule}{%
  {\color{colborder}\rule{\linewidth}{0.4pt}}%
}

%% ——— Intestazione sezione ————————————————————
\newcommand{\cvsection}[1]{%
  \vspace{0.7em}%
  {\small\bfseries\spaced[90]{\MakeUppercase{#1}}}%
  \vspace{0.15em}\newline
  \cvrule%
  \vspace{0.4em}\newline
}

%% ——— Voce opera: anno, titolo (italic), organico ——
\newcommand{\workrow}[3]{%
  {\footnotesize\color{coltextlight}\spaced[90]{#1}}\enspace
  {\small\itshape #2}\newline
  {\footnotesize\color{coltextlight}\hspace{2.2em}#3}%
  \vspace{0.1em}\newline
}

%% ——— Voce formazione / premio ————————————————
\newcommand{\cvrow}[3]{%
  {\small\bfseries #1}\hfill{\footnotesize\color{coltextlight}\spaced[90]{#2}}%
  \newline
  {\small\color{coltextlight}#3}%
  \vspace{0.55em}\newline
}

%% ===================================================
\begin{document}
\pagestyle{empty}

%% ——— INTESTAZIONE ———————————————————————————

\noindent\begin{minipage}[t]{0.60\linewidth}
  \vspace{0pt}
  {\fontsize{34}{38}\selectfont\spaced[30]{ALESSIO ROMA}}%
  \vspace{0.3em}\newline
  {\footnotesize\bfseries\color{coltextlight}\spaced[130]{MUSICISTA · COMPOSITORE}}%
  \vspace{0.55em}\newline
  {\color{colborder}\rule{0.97\linewidth}{0.4pt}}%
  \vspace{0.45em}\newline
  {\small\color{coltextlight}%
    Bari, 1993\enspace·\enspace
    \href{mailto:info@alessioroma.it}{info@alessioroma.it}\enspace·\enspace
    \href{https://alessioroma.it}{alessioroma.it}%
  }%
\end{minipage}%
\hfill%
\begin{minipage}[t]{0.31\linewidth}
  \vspace{0pt}
  \raggedleft
  \begingroup
    \setlength{\fboxsep}{0pt}%
    \setlength{\fboxrule}{1pt}%
    \color{coltextmain}%
    \fbox{%
      \includegraphics[width=3.8cm, height=3.8cm]%
        {assets/images/alessio-roma-smile-crop.jpg}%
    }%
  \endgroup
\end{minipage}

\vspace{0.9em}
\cvrule

%% ——— PROFILO ————————————————————————————————
\vspace{0.7em}
{\small Alessio Roma è un musicista e compositore nato a Bari nel 1993.
Nel suo lavoro integra elementi extra-musicali — salute mentale,
oppressione dei popoli, violenza sulle donne — attraverso le tecniche
spettraliste e la ricerca sulla fusione timbrica di più strumenti in
un unico timbro percepibile. Le sue opere sono eseguite in Italia e
all'estero in concerti, festival e rassegne di musica contemporanea.}

%% ——— FORMAZIONE ————————————————————————————
\cvsection{Formazione}

\cvrow{Conservatorio «Niccolò Piccinni», Bari}{2026}{%
  Diploma accademico di II livello in Composizione%
  \enspace·\enspace docente M°\,Biagio Putignano}

%% ——— MASTERCLASS ————————————————————————————
\cvsection{Masterclass internazionali}

{\small
Alessandro Melchiorre\enspace·\enspace
Ivan Fedele\enspace·\enspace
Giorgio Nottoli\enspace·\enspace
Marco Stroppa\enspace·\enspace
Robert HP Platz\enspace·\enspace
Roberto Paci Dalò\enspace·\enspace
Juha T.\,Koskinen\enspace·\enspace
Musikschule der Stadt Freising, Germania%
}

%% ——— PREMI ——————————————————————————————————
\cvsection{Premi e riconoscimenti}

{\small 2° Premio, Concorso di composizione «Note per una stella», Monteroni di Lecce (LE)}%
\newline
{\small\color{coltextlight}\hspace{0.5em}per \textit{Cinque piccole variazioni per tromba}}%
\vspace{0.55em}\newline

{\small Due prime esecuzioni assolute, eventi del centenario della\\
Società Italiana di Musica Contemporanea (SIMC)}%
\vspace{0.55em}\newline

%% ——— COMMISSIONI ————————————————————————————
\cvsection{Commissioni}

{\footnotesize\color{coltextlight}\spaced[90]{2026}}\enspace
{\small\textit{Ostinate Migrazioni} per quartetto d'archi — Traetta Opera Festival}%
\vspace{0.55em}\newline

{\footnotesize\color{coltextlight}\spaced[90]{2025}}\enspace
{\small\textit{Chiari Ombre} per violino e violoncello — Traetta Opera Festival}%
\vspace{0.55em}\newline

%% ——— COLLABORAZIONI ——————————————————————————
\cvsection{Collaborazioni}

{\small LEAP — Laboratorio ElettroAcustico Permanente, Roma}%
\newline
{\small\color{coltextlight}\hspace{0.5em}con Giuseppe Silvi (cofondatore)}

%% ——— DISCOGRAFIA ————————————————————————————
\cvsection{Discografia}

{\small Registrazioni disponibili su Spotify, Apple Music, YouTube Music e altre
piattaforme di streaming musicale.}

%% ===================================================
%% PAGINA 2 — CATALOGO DELLE OPERE
%% ===================================================
\newpage

\noindent
{\footnotesize\color{coltextlight}\spaced[100]{ALESSIO ROMA}}%
\hfill%
{\footnotesize\color{coltextlight}\spaced[100]{CATALOGO DELLE OPERE}}%
\vspace{0.3em}\newline
\cvrule

\setlength{\columnsep}{1.8em}
\begin{multicols}{2}

\cvsection{Musica strumentale}

\workrow{2026}{Ostinate Migrazioni}%
  {quartetto d'archi%
  \newline\hspace*{2.2em}{\itshape commissione Traetta Opera Festival}}

\workrow{2026}{Io sono come il mare}%
  {ensemble, elett. dal vivo, voce narrante — 45' ca.%
  \newline\hspace*{2.2em}testi di Antonella Diacono}

\workrow{2025}{Chiari Ombre}%
  {violino e violoncello — 6' ca.%
  \newline\hspace*{2.2em}{\itshape commissione Traetta Opera Festival}}

\workrow{2025}{Torsioni del reale}%
  {violino solo — 7' ca.}

\workrow{2025}{Fecero un deserto e lo chiamarono Pace}%
  {ensemble, elett. dal vivo, voce narrante — 15' ca.}

\workrow{2024}{Serenata per un attacco di panico III}%
  {pianoforte solo — 10' ca.}

\workrow{2024}{Costellazioni}%
  {orchestra — 9' ca.}

\workrow{2024}{Buttate in cielo}%
  {ensemble e narratore — 30' ca.%
  \newline\hspace*{2.2em}testi di Denise De Scisciolo}

\workrow{2023}{Paysage}%
  {organo storico e due performer — 8' ca.}

\workrow{2023}{Posters}%
  {clarinetto e pianoforte — 9' ca.}

\workrow{2023}{Serenata per un attacco di panico II}%
  {sassofono contralto solo — 8' ca.}

\workrow{2023}{Parla, inciampa, rifletti e splendi!}%
  {violino solo — 10' ca.}

\workrow{2023}{It's a long, long way to therapy}%
  {quartetto d'archi e elett. dal vivo — 8' ca.}

\workrow{2023{\addfontfeatures{LetterSpace=0}\;{\tiny(rev.\,2024)}}}{Serenata per un attacco di panico I}%
  {viola sola — 10' ca.}

\workrow{2022}{Fili d'aquilone}%
  {violino, violoncello e pianoforte — 10' ca.}

\workrow{2022}{Il Gioco Grande}%
  {sax baritono, percussioni e narratori — 15'–18' ca.}

\workrow{2022}{Boards}%
  {pianoforte solo — 7' ca.}

\workrow{2021}{Umarmung}%
  {duo di flauto dolce — 6' ca.}

\cvsection{Colonne sonore}

\workrow{2023}{Nu ffischia}%
  {cortometraggio di Minafra, Carella e Piccinni}

\workrow{2021}{Myosotys}%
  {cortometraggio di Minafra e Piccinni}

\workrow{2021}{Punto e a caso}%
  {cortometraggio di Minafra e Piccinni}

\end{multicols}

%% ——— Footer pagina 2 ————————————————————————
\vfill
\cvrule
\vspace{0.35em}
{\footnotesize\color{coltextlight}%
  \href{https://alessioroma.it}{alessioroma.it}%
  \enspace·\enspace
  \href{mailto:info@alessioroma.it}{info@alessioroma.it}%
  \enspace·\enspace
  Aggiornato: ottobre 2026%
}

\end{document}
"
%% ===================================================
\documentclass[a4paper, 10pt]{article}

\usepackage{fontspec}
\usepackage{xcolor}
\usepackage{geometry}
\usepackage{graphicx}
\usepackage{hyperref}
\usepackage{enumitem}
\usepackage{multicol}
\usepackage{setspace}

%% ——— Font ————————————————————————————————————
\setmainfont[
  Path           = assets/fonts/,
  UprightFont    = Champagne___Limousines.ttf,
  BoldFont       = Champagne___Limousines_Bold.ttf,
  ItalicFont     = Champagne___Limousines_Italic.ttf,
  BoldItalicFont = Champagne___Limousines_Bold_Italic.ttf
]{Champagne___Limousines.ttf}

%% ——— Colori dal sito ————————————————————————
\ifdefined\lightmode
  %% Versione chiara: stessi toni freddi del sito, scuriti per la stampa
  \definecolor{colbg}{HTML}{FFFFFF}
  \definecolor{coltextmain}{HTML}{1B2124}
  \definecolor{coltextlight}{HTML}{3F4C53}
  \definecolor{colaccent}{HTML}{2A5670}
  \definecolor{colborder}{HTML}{B9C4CA}
\else
  \definecolor{colbg}{HTML}{0D0F10}
  \definecolor{coltextmain}{HTML}{E8E0D0}
  \definecolor{coltextlight}{HTML}{7A8A90}
  \definecolor{colaccent}{HTML}{B8C8D4}
  \definecolor{colborder}{HTML}{2E3A42}
\fi

%% ——— Impostazioni pagina ————————————————————
\geometry{a4paper, left=2.2cm, right=2.2cm, top=2cm, bottom=2cm}
\pagecolor{colbg}
\color{coltextmain}
\setlength{\parindent}{0pt}
\setlength{\parskip}{0pt}

%% ——— Link ————————————————————————————————————
\hypersetup{
  colorlinks = true,
  urlcolor   = colaccent,
  linkcolor  = colaccent,
  pdfauthor  = {Alessio Roma},
  pdftitle   = {Alessio Roma — Curriculum Artistico}
}

%% ——— Helper: testo con spaziatura lettere ————
\newcommand{\spaced}[2][80]{%
  {\addfontfeatures{LetterSpace=#1}#2}%
}

%% ——— Linea divisoria ————————————————————————
\newcommand{\cvrule}{%
  {\color{colborder}\rule{\linewidth}{0.4pt}}%
}

%% ——— Intestazione sezione ————————————————————
\newcommand{\cvsection}[1]{%
  \vspace{0.7em}%
  {\small\bfseries\spaced[90]{\MakeUppercase{#1}}}%
  \vspace{0.15em}\newline
  \cvrule%
  \vspace{0.4em}\newline
}

%% ——— Voce opera: anno, titolo (italic), organico ——
\newcommand{\workrow}[3]{%
  {\footnotesize\color{coltextlight}\spaced[90]{#1}}\enspace
  {\small\itshape #2}\newline
  {\footnotesize\color{coltextlight}\hspace{2.2em}#3}%
  \vspace{0.1em}\newline
}

%% ——— Voce formazione / premio ————————————————
\newcommand{\cvrow}[3]{%
  {\small\bfseries #1}\hfill{\footnotesize\color{coltextlight}\spaced[90]{#2}}%
  \newline
  {\small\color{coltextlight}#3}%
  \vspace{0.55em}\newline
}

%% ===================================================
\begin{document}
\pagestyle{empty}

%% ——— INTESTAZIONE ———————————————————————————

\noindent\begin{minipage}[t]{0.60\linewidth}
  \vspace{0pt}
  {\fontsize{34}{38}\selectfont\spaced[30]{ALESSIO ROMA}}%
  \vspace{0.3em}\newline
  {\footnotesize\bfseries\color{coltextlight}\spaced[130]{MUSICISTA · COMPOSITORE}}%
  \vspace{0.55em}\newline
  {\color{colborder}\rule{0.97\linewidth}{0.4pt}}%
  \vspace{0.45em}\newline
  {\small\color{coltextlight}%
    Bari, 1993\enspace·\enspace
    \href{mailto:info@alessioroma.it}{info@alessioroma.it}\enspace·\enspace
    \href{https://alessioroma.it}{alessioroma.it}%
  }%
\end{minipage}%
\hfill%
\begin{minipage}[t]{0.31\linewidth}
  \vspace{0pt}
  \raggedleft
  \begingroup
    \setlength{\fboxsep}{0pt}%
    \setlength{\fboxrule}{1pt}%
    \color{coltextmain}%
    \fbox{%
      \includegraphics[width=3.8cm, height=3.8cm]%
        {assets/images/alessio-roma-smile-crop.jpg}%
    }%
  \endgroup
\end{minipage}

\vspace{0.9em}
\cvrule

%% ——— PROFILO ————————————————————————————————
\vspace{0.7em}
{\small Alessio Roma è un musicista e compositore nato a Bari nel 1993.
Nel suo lavoro integra elementi extra-musicali — salute mentale,
oppressione dei popoli, violenza sulle donne — attraverso le tecniche
spettraliste e la ricerca sulla fusione timbrica di più strumenti in
un unico timbro percepibile. Le sue opere sono eseguite in Italia e
all'estero in concerti, festival e rassegne di musica contemporanea.}

%% ——— FORMAZIONE ————————————————————————————
\cvsection{Formazione}

\cvrow{Conservatorio «Niccolò Piccinni», Bari}{2026}{%
  Diploma accademico di II livello in Composizione%
  \enspace·\enspace docente M°\,Biagio Putignano}

%% ——— MASTERCLASS ————————————————————————————
\cvsection{Masterclass internazionali}

{\small
Alessandro Melchiorre\enspace·\enspace
Ivan Fedele\enspace·\enspace
Giorgio Nottoli\enspace·\enspace
Marco Stroppa\enspace·\enspace
Robert HP Platz\enspace·\enspace
Roberto Paci Dalò\enspace·\enspace
Juha T.\,Koskinen\enspace·\enspace
Musikschule der Stadt Freising, Germania%
}

%% ——— PREMI ——————————————————————————————————
\cvsection{Premi e riconoscimenti}

{\small 2° Premio, Concorso di composizione «Note per una stella», Monteroni di Lecce (LE)}%
\newline
{\small\color{coltextlight}\hspace{0.5em}per \textit{Cinque piccole variazioni per tromba}}%
\vspace{0.55em}\newline

{\small Due prime esecuzioni assolute, eventi del centenario della\\
Società Italiana di Musica Contemporanea (SIMC)}%
\vspace{0.55em}\newline

%% ——— COMMISSIONI ————————————————————————————
\cvsection{Commissioni}

{\footnotesize\color{coltextlight}\spaced[90]{2026}}\enspace
{\small\textit{Ostinate Migrazioni} per quartetto d'archi — Traetta Opera Festival}%
\vspace{0.55em}\newline

{\footnotesize\color{coltextlight}\spaced[90]{2025}}\enspace
{\small\textit{Chiari Ombre} per violino e violoncello — Traetta Opera Festival}%
\vspace{0.55em}\newline

%% ——— COLLABORAZIONI ——————————————————————————
\cvsection{Collaborazioni}

{\small LEAP — Laboratorio ElettroAcustico Permanente, Roma}%
\newline
{\small\color{coltextlight}\hspace{0.5em}con Giuseppe Silvi (cofondatore)}

%% ——— DISCOGRAFIA ————————————————————————————
\cvsection{Discografia}

{\small Registrazioni disponibili su Spotify, Apple Music, YouTube Music e altre
piattaforme di streaming musicale.}

%% ===================================================
%% PAGINA 2 — CATALOGO DELLE OPERE
%% ===================================================
\newpage

\noindent
{\footnotesize\color{coltextlight}\spaced[100]{ALESSIO ROMA}}%
\hfill%
{\footnotesize\color{coltextlight}\spaced[100]{CATALOGO DELLE OPERE}}%
\vspace{0.3em}\newline
\cvrule

\setlength{\columnsep}{1.8em}
\begin{multicols}{2}

\cvsection{Musica strumentale}

\workrow{2026}{Ostinate Migrazioni}%
  {quartetto d'archi%
  \newline\hspace*{2.2em}{\itshape commissione Traetta Opera Festival}}

\workrow{2026}{Io sono come il mare}%
  {ensemble, elett. dal vivo, voce narrante — 45' ca.%
  \newline\hspace*{2.2em}testi di Antonella Diacono}

\workrow{2025}{Chiari Ombre}%
  {violino e violoncello — 6' ca.%
  \newline\hspace*{2.2em}{\itshape commissione Traetta Opera Festival}}

\workrow{2025}{Torsioni del reale}%
  {violino solo — 7' ca.}

\workrow{2025}{Fecero un deserto e lo chiamarono Pace}%
  {ensemble, elett. dal vivo, voce narrante — 15' ca.}

\workrow{2024}{Serenata per un attacco di panico III}%
  {pianoforte solo — 10' ca.}

\workrow{2024}{Costellazioni}%
  {orchestra — 9' ca.}

\workrow{2024}{Buttate in cielo}%
  {ensemble e narratore — 30' ca.%
  \newline\hspace*{2.2em}testi di Denise De Scisciolo}

\workrow{2023}{Paysage}%
  {organo storico e due performer — 8' ca.}

\workrow{2023}{Posters}%
  {clarinetto e pianoforte — 9' ca.}

\workrow{2023}{Serenata per un attacco di panico II}%
  {sassofono contralto solo — 8' ca.}

\workrow{2023}{Parla, inciampa, rifletti e splendi!}%
  {violino solo — 10' ca.}

\workrow{2023}{It's a long, long way to therapy}%
  {quartetto d'archi e elett. dal vivo — 8' ca.}

\workrow{2023{\addfontfeatures{LetterSpace=0}\;{\tiny(rev.\,2024)}}}{Serenata per un attacco di panico I}%
  {viola sola — 10' ca.}

\workrow{2022}{Fili d'aquilone}%
  {violino, violoncello e pianoforte — 10' ca.}

\workrow{2022}{Il Gioco Grande}%
  {sax baritono, percussioni e narratori — 15'–18' ca.}

\workrow{2022}{Boards}%
  {pianoforte solo — 7' ca.}

\workrow{2021}{Umarmung}%
  {duo di flauto dolce — 6' ca.}

\cvsection{Colonne sonore}

\workrow{2023}{Nu ffischia}%
  {cortometraggio di Minafra, Carella e Piccinni}

\workrow{2021}{Myosotys}%
  {cortometraggio di Minafra e Piccinni}

\workrow{2021}{Punto e a caso}%
  {cortometraggio di Minafra e Piccinni}

\end{multicols}

%% ——— Footer pagina 2 ————————————————————————
\vfill
\cvrule
\vspace{0.35em}
{\footnotesize\color{coltextlight}%
  \href{https://alessioroma.it}{alessioroma.it}%
  \enspace·\enspace
  \href{mailto:info@alessioroma.it}{info@alessioroma.it}%
  \enspace·\enspace
  Aggiornato: ottobre 2026%
}

\end{document}
